\chapter{Multiband Aperture and Polarimetry}

The aperture is intentionally hybrid rather than uniform across the spectrum. One external single-polarization VHF Yagi and one external single-polarization UHF Yagi share a boom aligned with the outward face normal: the VHF antenna uses larger elements and the UHF antenna uses smaller elements. Their transverse directions are mutually orthogonal at 45° and 135° in the local tangent basis. These mechanical orientations provide diversity across different bands, not two polarization components of one signal. Smaller L/S/C, X/Ku and K/Ka antennas or tiles occupy the face layers and local mounts. HF uses loops, crossed dipoles and platform-mode diversity. Each antenna family has its own filtering, conversion and sampling strategy, while the face preserves a common timing, calibration and fusion interface.

The first demonstrator does not claim continuous spectral coverage. The 323--470 MHz and 860--960 MHz intervals are deliberately unassigned. A dedicated 960--1215 MHz aviation aperture or array supplies two simultaneous independent paths: a 978 MHz UAT preselector, limiter/LNA, coherent ADC and FPGA channel; and a 1090 MHz 1090ES preselector, limiter/LNA, coherent ADC and FPGA channel. Neither path is assigned to the 470--860 MHz UHF Yagi, and the aviation aperture does not imply continuous monitoring across its full tuning range.

Polarimetric synthesis is permitted only in a same-band module that provides two simultaneous, coherent orthogonal ports at the same frequency (\texttt{POL-001}). In a local $x$--$y$ basis:
\[
E_{\mathrm{RHCP}}=\frac{E_x-jE_y}{\sqrt{2}},\qquad E_{\mathrm{LHCP}}=\frac{E_x+jE_y}{\sqrt{2}}.
\]
The equations do not combine the VHF and UHF Yagis. The first VHF/UHF demonstrator treats them as independent single-polarization channels (\texttt{POL-002}); full polarimetry requires a later dual-port or same-band crossed module. For a curved surface, valid local bases must be rotated to a common coordinate system before fusion. Polarimetric calibration is a function of frequency, angle, temperature, amplitude balance, phase and mutual coupling.

A separate future four-channel mast, distinct from the illustrated combined VHF/UHF pair, implements VHF-X/VHF-Y at 0°/90° and UHF-A/UHF-B at 45°/135°, with four electrically independent booms and RF chains (\texttt{POL-004}--\texttt{POL-006}). For each band $b$, calibration follows:
\[
\widehat{\mathbf{E}}_b(f,\theta,\phi,T)=
\mathbf{C}_b^{-1}(f,\theta,\phi,T)
\left[\mathbf{v}_b-\widehat{\mathbf{n}}_b\right],
\]
where $\mathbf{C}_b$ is the measured complex $2\times2$ Jones matrix containing embedded patterns, gain/phase imbalance, delays, mutual coupling, mast and platform effects. A regularized pseudoinverse replaces the direct inverse where conditioning is poor. The UHF estimate is then rotated by $\psi=45^\circ$ into the common global basis before polarization products are formed.

\section{Engineering solution}
The crossed-band Yagi assembly is a proposed mechanical integration solution. The larger Yagi covers VHF and the smaller Yagi covers UHF; each remains independently matched, digitized and calibrated. Their outputs may be compared or fused as different-band observations, but no relative phase operation produces RHCP, LHCP, Jones or Stokes quantities from this pair. Novelty and performance are not claimed until prior-art review and measurement.

For wider VHF discovery coverage, an 11-element 72--320 MHz LPDA with a 2 m maximum element is retained as a simulation candidate (\texttt{ANT-001}--\texttt{ANT-005}). It is not yet a replacement for the illustrated VHF Yagi: the approximately 2 m boom, face-edge clearance, wind load, feed transition, gain and interaction with the UHF antenna require validation.

\begin{figure}[H]
\centering
\includegraphics[width=0.95\textwidth]{en/fig04_particionamento_espectro.png}
\caption{Spectral partition on a logarithmic horizontal scale. Each coloured bar is one proposed band family; the blank intervals between the VHF, UHF and dedicated aviation paths are deliberate first-demonstrator gaps.}
\label{fig:fig04_particionamento_espectro}
\end{figure}

\begin{figure}[H]
\centering
\includegraphics[width=0.95\textwidth]{en/fig05_polarimetria.png}
\caption{Top flow: independently digitized VHF and UHF channels terminate in a crossed-out cross-band polarimetric result. Bottom flow: two coherent orthogonal ports in the same band feed calibrated Jones, Stokes or circular-polarization products.}
\label{fig:fig05_polarimetria}
\end{figure}
